% ============================================================
%  example_hard.tex  ——  挑战示例：真实论文里的"暗坑"
%  用法：math_guard.exe example_hard.tex
%
%  这里的每一处都来自真实论文（arXiv / CTAN / 用户稿件），
%  而且是"编译 0 报错、日志 0 错误、PDF 印出来是错的"那种。
% ============================================================
\documentclass{article}
\usepackage{amsmath,amssymb}
\usepackage{hyperref}

% 自定义宏：工具会自动展开
\newcommand{\hf}{\frac{1}{2}}
\newcommand{\RR}{\mathbb{R}}
\newcommand{\eqd}{\overset{d}{=}}
\def\sls#1#2{(\frac{#1}{#2})}

\begin{document}

\section{真实案例：55 页论文里的空分母}

% 这处的原文是：
%   {}_{2}F_{1}\Bigl(\hf,\hf;1;\frac{-(e^{-x}-1)^{2}}{}\Bigr)
% xelatex 编译 exit=0，日志 0 错误，PDF 第 38 页印成一道空横线。
% 只有肉眼盯着看才看得出来。
\[
{}_{2}F_{1}\Bigl(\hf,\hf;1;4e^{x}(1-e^{x})\Bigr)=(2e^{x}-1)^{-1}
{}_{2}F_{1}\Bigl(\hf,\hf;1;\frac{-(e^{-x}-1)^{2}}{(e^{x}-1)^{2}}\Bigr)
\]
% 上面这个已经修好了；把分母改成空的就是原文的状态。

\section{自定义宏展开}

% \def 定义的带参宏：展开后是 (1+\sqrt{-7})/2
\[ \sls{1+\sqrt{-7}}{2}2^{24} = 2^{47}+2^{23}\sqrt{-7} \]

% 无参宏
\[ \hf+\hf=1 \]

\section{各种"看起来像错、其实没错"的写法}

% 分式等于小数（浮点精度陷阱）
\[ \frac{1}{200}=0.005 \]

% 百分比
\[ \frac{26}{500}=5.2\% \]

% 括号幂（曾被误判）
\[ 5(5^{2}-1)=120 \]

% 幂是表达式
\[ 2^{4-1}=8 \]

% 方括号做分组
\[ \frac{8-(11^{2}+3^{2})/21}{4}=\frac{19}{42} \]

% 有限小数舍入
\[ \frac{0.5}{(0.25+1)^{1/2}}=0.447 \]

% 浮点累积误差
\[ 1\times 0.175^{2}-0.5\times 0.195^{2}=0.0116125 \]

\section{真错误（这些确实要报）}

% 链式等式里第二段算错
\[ r(4)=16-6=11 \]

% 整数与分数混算
\[ 1=\frac{1}{2} \]

% 不等式方向错
\[ 40\ge 6\cdot 20 \]

\section{环境与引用}

\begin{equation}\label{eq:ok}
  E = mc^{2}
\end{equation}

引用 \eqref{eq:ok} 是正常的；引用 \eqref{eq:missing} 会报未定义。

% 环境不配对：begin 是 equation，end 是 align
\begin{equation}
  y = x + 1
\end{align}

\section{无法验证（工具会承认）}

% 概率论记号
\[ X_{n}\eqd N(0,1) \]

% 矩阵
\[ A=\begin{pmatrix} 1 & 2 \\ 3 & 4 \end{pmatrix} \]

% 纯符号恒等式
\[ (a+b)^{2}=a^{2}+2ab+b^{2} \]

\end{document}
